\chapter{আলোচনা}\label{ch:discussion}

\section{গবেষণা প্রশ্নের উত্তর}

\paragraph{RQ1: ডেটা অ্যানালিটিক্স কীভাবে ব্যবসায়িক সিদ্ধান্ত গ্রহণ উন্নত করতে পারে?}
কাঠামোটি অ্যানালিটিক্সের তিনটি স্তরেই উন্নতি দেখায়। \emph{বর্ণনামূলক} স্তরে এটি হাতে-করা স্প্রেডশিট একত্রীকরণের বদলে যাচাইকৃত ও
মিলানো KPI দেয়, যা একটি কমান্ডে হালনাগাদ হয়। \emph{রোগনির্ণয়মূলক} স্তরে এটি এমন সম্পর্ক উন্মোচন করে যা সাধারণ প্রতিবেদনে চোখ
এড়িয়ে যেত: প্রবৃদ্ধি মূলত মূল্য ও মিশ্রণ-চালিত; ডিসকাউন্ট বড় অর্ডারের সঙ্গে যুক্ত নয়; অর্ডারের আকার গ্রাহকের ধরন নয়, চ্যানেলের উপর নির্ভর
করে। \emph{পূর্বাভাসমূলক} স্তরে এটি পরিমাপিত অনিশ্চয়তাসহ পণ্য-স্তরের পূর্বাভাস দেয়। এর প্রতিটি একটি সুনির্দিষ্ট সিদ্ধান্ত বদলায়: ডিসকাউন্ট
বাজেট, বিক্রয়-বাহিনী বণ্টন, উৎপাদন পরিকল্পনা ও লজিস্টিক সক্ষমতা।

\paragraph{RQ2: কোন প্রিপ্রসেসিং কৌশল সবচেয়ে কার্যকর?}
তিনটি কৌশল সবচেয়ে মূল্যবান প্রমাণিত হয়েছে। প্রথমত, হিসাবগত অভেদ থেকে অনুপস্থিত মান \emph{নির্ণয়}, যা হুবহু পুনরুদ্ধার দেয় এবং
পরিসংখ্যানিক imputation-এর আগে সবসময় চেষ্টা করা উচিত। দ্বিতীয়ত, \emph{robust, দলভিত্তিক outlier শনাক্তকরণ} (log স্কেলে প্রতি-পণ্য
IQR), যা বিষম বণ্টন ও ভিন্ন পণ্য-স্কেলকে সম্মান করে; অস্বাভাবিক মান মুছে না ফেলে চিহ্নিত করার নীতির সঙ্গে মিলে এটি প্রকৃত বাল্ক অর্ডারগুলো
রক্ষা করে। তৃতীয়ত, \emph{leakage-মুক্ত ফিচার ইঞ্জিনিয়ারিং}, যেখানে rolling-এর আগে ইতিহাস-ফিচার shift করা হয় এবং প্রতিটি শেখা প্যারামিটার
কেবল প্রশিক্ষণ ডেটায় fit করা হয়। ব্যবসায়িক চলকগুলো অত্যন্ত বিষম হওয়ায় সর্বত্র normal-তত্ত্বের পদ্ধতির চেয়ে robust পদ্ধতি শ্রেয় ছিল।

\paragraph{RQ3: কোন অ্যালগরিদম সবচেয়ে ভালো?}
সব সমস্যায় একটি অ্যালগরিদম সেরা নয়, এবং ভালো মডেলগুলোর ব্যবধান সামান্য। এক মাস অগ্রিম পণ্য-চাহিদায় \BestReg সবচেয়ে নির্ভুল ছিল;
বহু-ধাপ পরিস্থিতিতে \BestRec সবচেয়ে দৃঢ়; কোম্পানি-পর্যায়ে ধ্রুপদী Holt--Winters মডেল সব মেশিন লার্নিং বিকল্পকে হারিয়েছে; আর
ক্লাসিফিকেশনে logistic regression tree ensemble-এর সমান। এই ধারা \citet{makridakis2018}-এর সঙ্গে মেলে: অ্যালগরিদম বাছাইয়ের চেয়ে
ভালো ফিচার, সঠিক যাচাই-নকশা ও বেঞ্চমার্ক বেশি গুরুত্বপূর্ণ। ব্যবহারিকভাবে সবচেয়ে উপযোগী একটি সমন্বিত পদ্ধতি: পরিসংখ্যানিক মোট, মেশিন
লার্নিং দিয়ে শেখা পণ্য-মিশ্রণে মিলানো।

\paragraph{RQ4: ঐতিহাসিক ডেটায় কী প্যাটার্ন আছে?}
ডিসেম্বরে শীর্ষে পৌঁছানো শক্তিশালী বছর-শেষের মৌসুমিতা; স্থির কিন্তু মূল্য-চালিত প্রবৃদ্ধি ও মূল্য-অস্থিতিস্থাপক চাহিদা; ঢাকা ও চট্টগ্রামে
ভৌগোলিকভাবে এবং প্রায় দশটি পণ্যে পোর্টফোলিও-গতভাবে কেন্দ্রীভূত আয়; চ্যানেলভেদে অর্ডারের আকারে বড় পার্থক্য; ডিসকাউন্টের কোনো শনাক্তযোগ্য
পরিমাণ-প্রভাব নেই; এবং সপ্তাহের দিনের প্রভাব নেই। বিভাজন যোগ করে একটি পোর্টফোলিও-দৃষ্টিভঙ্গি (প্রিমিয়াম ও গিফট প্যাক, কোর ফ্যামিলি
প্যাক, ভ্যালু সিঙ্গেল-সার্ভ ও নিশ লাইন) এবং একটি বাজার-দৃষ্টিভঙ্গি (core retail, bulk distribution, online ও আঞ্চলিক, উদীয়মান
উচ্চ-প্রবৃদ্ধি)।

\paragraph{RQ5: ভবিষ্যৎ কতটা নির্ভুলভাবে পূর্বাভাস করা যায়?}
প্রতিটি পণ্যের মাসিক চাহিদা এক মাস অগ্রিম প্রায় \BestRegWAPE\% এবং 21 মাস পর্যন্ত প্রায় \BestRecWAPE\% weighted ত্রুটিতে পূর্বাভাস করা
যায়; কোম্পানি-স্তরের মাসিক ইউনিট প্রায় \HWMAPE\% ত্রুটিতে। একটি নির্দিষ্ট অর্ডার বড় হবে কি না তা কেবল মাঝারি মাত্রায় পূর্বাভাসযোগ্য
(ROC-AUC \ClsAUC); সীমাবদ্ধতা মডেলে নয়, প্রাপ্ত তথ্যে।

\paragraph{RQ6: ব্যবস্থাপকদের কাছে ফলাফল কীভাবে উপস্থাপন করা উচিত?}
ব্যবস্থাপনা প্রশ্নগুলোকে ঘিরে সাজানো একটি ইন্টারঅ্যাকটিভ ড্যাশবোর্ডের মাধ্যমে, যেখানে সর্বজনীন ফিল্টার, KPI কার্ড ও সরল সাজানো চার্ট থাকবে,
পূর্বাভাসে সবসময় অনিশ্চয়তা দেখানো হবে, এবং প্রতিটি অন্তর্দৃষ্টি একটি করণীয়সহ সহজ ভাষায় বলা হবে। বিনামূল্যের প্ল্যাটফর্মে ড্যাশবোর্ড ও
সারসংক্ষেপ ওয়েবসাইট প্রকাশ করলে খরচ ও ইনস্টলেশনের বাধা দূর হয়।

\paragraph{RQ7: ঐতিহাসিক প্রতিবেদনের তুলনায় পূর্বাভাসমূলক অ্যানালিটিক্স কতটা যোগ করে?}
একটি ঐতিহাসিক প্রতিবেদন বড়জোর গত বছরের প্যাটার্ন পুনরাবৃত্তি করতে পারে, যা seasonal-naive নিয়ম। মেশিন লার্নিং পূর্বাভাস এর ত্রুটি
\GainVsSNaive\% কমায়। অন্যভাবে বললে, গত বছরের হিসাব দিয়ে পরিকল্পনাকারী প্রতি পণ্য-মাসে পরিমাণের প্রায় \SNaiveWAPE\% ভুল করতেন, যেখানে
এই কাঠামোয় তা \BestRegWAPE\%। এর উপর, পূর্বাভাস-ব্যবধান এমন তথ্য যোগ করে যা ঐতিহাসিক প্রতিবেদন একেবারেই দিতে পারে না: কতটা
নিরাপত্তা-মজুদ রাখতে হবে।

\section{ব্যবস্থাপনাগত তাৎপর্য}
\begin{itemize}
  \item \textbf{মূল্য ও প্রমোশন।} চাহিদা অস্থিতিস্থাপক এবং ডিসকাউন্ট অর্ডারের আকার বাড়ায় না। এই সময়কালে প্রায় \DiscountCost মিলিয়ন টাকার
        ডিসকাউন্ট বাজেট পুনর্বিবেচনা করা উচিত, এবং নিয়ন্ত্রিত পরীক্ষার (A/B প্রমোশন) মাধ্যমে যাচাই করা উচিত লক্ষ্যভিত্তিক প্রমোশন লাভজনক কি না।
  \item \textbf{পরিমাণ-প্রবৃদ্ধি।} প্রবৃদ্ধি মূল্য-নির্ভর হওয়ায় পরিমাণকে আলাদা KPI হিসেবে অনুসরণ করা উচিত, এবং কম-প্রবেশিত অঞ্চল ও
        E-commerce-এ প্রবৃদ্ধির উদ্যোগ নেওয়া উচিত।
  \item \textbf{মৌসুমি পরিকল্পনা।} ডিসেম্বরের শীর্ষের জন্য অক্টোবর থেকে উৎপাদন ও ডিস্ট্রিবিউটর-মজুদ বাড়াতে হবে; ফেব্রুয়ারি রক্ষণাবেক্ষণ ও
        প্রমোশনের স্বাভাবিক সময়।
  \item \textbf{পোর্টফোলিও।} A-শ্রেণির পণ্যগুলো সর্বোচ্চ সেবা-স্তরের যোগ্য; C-শ্রেণির পণ্যগুলো পর্যালোচনার প্রার্থী।
  \item \textbf{লজিস্টিকস।} বড় অর্ডার ডিস্ট্রিবিউটর চ্যানেল ও নির্দিষ্ট অঞ্চলে কেন্দ্রীভূত, তাই যানবাহন ও মজুদ বণ্টন সেই স্তরে পরিকল্পনা করা যায়।
  \item \textbf{নিয়ন্ত্রণ।} দৈনিক অ্যানোমালি যাচাই ডেটার মান উন্নত করে এবং পরিচালনাগত সমস্যা আগেভাগে ধরে।
\end{itemize}

\section{সীমাবদ্ধতা}\label{sec:limitations}
\begin{itemize}
  \item \textbf{ডেটার উৎস ও বাস্তবতা।} ডেটাসেটটি অত্যন্ত পরিষ্কার: কোনো অনুপস্থিত মান নেই, হিসাবগত অভেদ হুবহু মেলে, এবং অর্ডারের আকার
        গ্রাহকের ধরন, ডিসকাউন্ট ও পণ্য থেকে প্রায় স্বাধীন। বাস্তব পরিচালনাগত ডেটা সাধারণত এর চেয়ে অগোছালো। তাই পাইপলাইনটি ইচ্ছাকৃতভাবে
        নষ্ট করা অনুলিপিতে পরীক্ষা করা হয়েছে, তবে গ্রাহক-আচরণ সম্পর্কিত সিদ্ধান্তে পদক্ষেপ নেওয়ার আগে কোম্পানির আরও ডেটায় যাচাই করা উচিত।
        ডেটাসেটের সঠিক উৎস ও সংগ্রহ-প্রক্রিয়া এই প্রতিবেদনের সঙ্গে নথিভুক্ত করা উচিত।
  \item \textbf{অনুপস্থিত চালিকাশক্তি।} বিপণন-ব্যয়, প্রতিযোগীর মূল্য, মজুদ-ঘাটতি, আবহাওয়া বা ছুটির কোনো ডেটা পাওয়া যায়নি; এগুলো যুক্ত হলে
        পূর্বাভাস ও ক্লাসিফিকেশন উভয়ই সম্ভবত উন্নত হবে।
  \item \textbf{পর্যবেক্ষণভিত্তিক স্থিতিস্থাপকতা।} স্থিতিস্থাপকতার অনুমান সময়ের সঙ্গে তালিকা-মূল্যের পরিবর্তনের উপর নির্ভরশীল, এলোমেলো পরীক্ষার
        উপর নয়; মূল্য সামান্য পরিমাণে বদলেছে, তাই পর্যবেক্ষিত সীমার বাইরে স্থিতিস্থাপকতা অজানা।
  \item \textbf{ভবিষ্যৎ পরিস্থিতি।} সব পূর্বাভাস ধরে নেয় অতীতের প্যাটার্ন চলবে; কাঠামোগত পরিবর্তন (যেমন নতুন প্রতিযোগী বা অর্থনৈতিক ধাক্কা)
        নির্ভুলতা কমাবে।
  \item \textbf{মূল্যায়নের পরিধি।} ড্যাশবোর্ডটি এখনো প্রকৃত ব্যবস্থাপকদের দিয়ে মূল্যায়িত হয়নি; আর প্যানেলটি (52টি সিরিজ, 81 মাস) এত ছোট যে
        deep-learning মডেল ব্যবহারের যৌক্তিকতা নেই, তাই সেগুলো পরীক্ষা করা হয়নি।
\end{itemize}

\section{পদ্ধতি নিয়ে প্রতিফলন}
সবচেয়ে গুরুত্বপূর্ণ পদ্ধতিগত সিদ্ধান্ত ছিল বেঞ্চমার্কসহ কঠোর সময়-ক্রমিক মূল্যায়ন; এলোমেলো cross-validation চমৎকার কিন্তু বিভ্রান্তিকর
নির্ভুলতার সংখ্যা দিত। কোড থেকে প্রতিটি উল্লিখিত সংখ্যা তৈরি করাও মূল্যবান প্রমাণিত হয়েছে: কাজ এগোনোর সঙ্গে বিশ্লেষণ বহুবার পুনরায় চালানো
হয়েছে, এবং প্রতিবেদন সবসময় সামঞ্জস্যপূর্ণ থেকেছে। এ ধরনের ভবিষ্যৎ প্রকল্পে দুটি চর্চাই সুপারিশযোগ্য।
